%! Author = chtompki
%! Date = 1/14/26
% Example LaTeX document for GP111 - note % sign indicates a comment
\documentclass[12pt]{amsart}
% Default margins are too wide all the way around. I reset them here
\setlength{\hoffset}{-.75in}
\setlength{\textwidth}{6.5in}
\setlength{\voffset}{-.5in}
\setlength{\textheight}{9.0in}
\usepackage{amsmath}
\usepackage{amssymb}
\usepackage{amsthm}
\usepackage{pgfplots}
\usepackage{enumerate}
\usepackage{tikz}
\usetikzlibrary{shadows,arrows,arrows.meta,shapes,positioning,calc}% Required for the Circle and Triangle arrow tips
\usepackage{setspace}
\usepackage{siunitx}
\usepackage{enumitem}

\theoremstyle{conjecture}
\newtheorem{conjecture}{Conjecture}

\theoremstyle{definition}
\newtheorem{definition}{Definition}

\theoremstyle{question}
\newtheorem{question}{Question}

\theoremstyle{theorem}
\newtheorem{theorem}{Theorem}

\newenvironment{solution}
{\begin{proof}[Solution]}
{\end{proof}}

\title{Graph Theory:\\Homework 6\\ Book Section 3.1 Part II.}
\author{Rob Tompkins}
\date{\today}


\begin{document}

    \maketitle
    \noindent\textbf{1.} Prove or disprove: Every tree has at most one perfect matching.
    (Hint: If a tree had two different perfect matchings, what could you say about the
    edges that belong to exactly one of them?)

    \begin{proof}
    The statement is true. Suppose, for a contradiction, that a tree \(T\)
    has distinct perfect matchings \(M_1\) and \(M_2\). Consider the spanning
    subgraph with edge set
    \[
        M_1\mathbin{\triangle}M_2
        =(M_1\setminus M_2)\cup(M_2\setminus M_1).
    \]
    At each vertex, the two perfect matchings either use the same edge,
    giving degree zero in this subgraph, or use different edges, giving
    degree two. Since the matchings are distinct, this subgraph has a
    component containing an edge. Every vertex of that component has
    degree two, so the component is a cycle. This contradicts the fact
    that a tree contains no cycle. Thus every tree has at most one
    perfect matching.
    \end{proof}

    \noindent\textbf{2.} Prove that every maximal matching in a graph \(G\) has at least
    \(\alpha'(G)/2\) edges.
    (Hint: Compare a maximal matching $M$ with a maximum matching $M^*$.
    What does maximality of $M$ tell you about each edge of the maximum matching?)

    \begin{proof}
    Let \(M\) be a maximal matching and let \(M^*\) be a maximum
    matching in \(G\). Every edge of \(M^*\) has at least one endpoint
    saturated by \(M\); otherwise, that edge could be added to \(M\),
    contradicting maximality.

    For each edge of \(M^*\), choose an edge of \(M\) sharing an endpoint
    with it. Each edge of \(M\) can be chosen by at most two edges of
    \(M^*\), since it has two endpoints and at most one edge of the
    matching \(M^*\) is incident with each endpoint. Consequently,
    \[
        \alpha'(G)=|M^*|\leq 2|M|,
    \]
    and hence \(|M|\geq \alpha'(G)/2\).
    \end{proof}

    \noindent\textbf{3.} Prove that if \(G\) is an \(X,Y\)-bigraph such that
    \[
        |N(S)| > |S|
        \qquad \text{whenever } \varnothing \neq S \subsetneq X,
    \]
    then every edge of \(G\) belongs to some matching that saturates \(X\).
    (Hint: Given an edge \(xy\), consider the graph
    \(G' = G-x-y\). What would you like to show about \(G'\), and how
    does deleting \(y\) affect neighborhoods?)

    \begin{proof}
    Fix any edge \(xy\in E(G)\), where \(x\in X\) and \(y\in Y\).
    Let \(G'=G-x-y\), with partite sets \(X'=X\setminus\{x\}\)
    and \(Y'=Y\setminus\{y\}\). We verify Hall's condition for \(X'\).
    If \(S\subseteq X'\) is nonempty, then \(S\) is a proper subset
    of \(X\), so the hypothesis and integrality give
    \(|N_G(S)|\geq |S|+1\). Deleting \(x\) does not affect the
    neighbors of \(S\), and deleting \(y\) removes at most one neighbor.
    Therefore,
    \[
        |N_{G'}(S)|=|N_G(S)\setminus\{y\}|
        \geq |N_G(S)|-1\geq |S|.
    \]
    The inequality also holds for \(S=\varnothing\). By Hall's Theorem,
    \(G'\) has a matching \(M'\) saturating \(X'\). Since neither
    \(x\) nor \(y\) belongs to \(G'\), the set \(M'\cup\{xy\}\)
    is a matching in \(G\) that saturates \(X\) and contains \(xy\).
    As \(xy\) was arbitrary, the conclusion follows. If \(G\) has
    no edges, the assertion about every edge is vacuous.
    \end{proof}

    \noindent\textbf{4.} Use the K\"onig--Egerv\'ary Theorem to prove that every bipartite
    graph \(G\) has a matching of size at least \(e(G)/\Delta(G)\).
    Use this to conclude that every subgraph of \(K_{n,n}\) with more
    than \((k-1)n\) edges has a matching of size at least \(k\).
    (Hint: K\"onig--Egerv\'ary lets you bound a different parameter
    instead of \(\alpha'(G)\). How many edges can a single vertex
    take care of?)

    \begin{proof}
    Assume \(\Delta(G)>0\), as required for the stated ratio to be
    defined. Let \(Q\) be a minimum vertex cover of \(G\). By the
    K\"onig--Egerv\'ary Theorem, \(|Q|=\alpha'(G)\). Since every
    edge has at least one endpoint in \(Q\),
    \[
        e(G)\leq \sum_{v\in Q}d_G(v)
        \leq |Q|\Delta(G)=\alpha'(G)\Delta(G).
    \]
    Dividing by \(\Delta(G)\) yields
    \(\alpha'(G)\geq e(G)/\Delta(G)\).
    For an edgeless graph, the empty matching has size zero; the
    division-free inequality above still holds.

    Now let \(G\) be a subgraph of \(K_{n,n}\) with
    \(e(G)>(k-1)n\), where \(n\) and \(k\) are positive integers.
    Then \(G\) has at least one edge and \(\Delta(G)\leq n\), so
    \[
        \alpha'(G)\geq \frac{e(G)}{\Delta(G)}
        \geq \frac{e(G)}{n}>k-1.
    \]
    Since \(\alpha'(G)\) is an integer, \(\alpha'(G)\geq k\).
    \end{proof}

    \noindent\textbf{5.} Let \(G\) be an \(X,Y\)-bigraph with vertex cover \(Q\), and let
    \(H\) be the induced \((X-Q),(Q\cap Y)\)-sub-bigraph of \(G\).
    Prove that if no matching in \(H\) saturates \(Q\cap Y=Q-X\),
    then \(G\) has a vertex cover that is smaller than \(Q\).
    (Hint: What does Hall's Theorem tell you when no matching in
    \(H\) saturates \(Q-X\)? Use that to decide which vertices of \(Q\)
    you might remove and replace to obtain a smaller set, and explain
    why the new set is a vertex cover.)

    \begin{proof}
    Since no matching in \(H\) saturates \(Q\cap Y\), Hall's
    Theorem gives a set \(S\subseteq Q\cap Y\) such that
    \[
        |N_H(S)|<|S|.
    \]
    Define
    \[
        Q'=(Q\setminus S)\cup N_H(S).
    \]
    Because \(N_H(S)\subseteq X\setminus Q\), the two sets in this
    union are disjoint, and
    \[
        |Q'|=|Q|-|S|+|N_H(S)|<|Q|.
    \]
    It remains to show that \(Q'\) is a vertex cover of \(G\).
    Let \(uv\) be any edge, with \(u\in X\) and \(v\in Y\).
    If \(v\notin S\), at least one endpoint lies in \(Q\), and
    neither endpoint lies in \(S\); thus an endpoint remains in
    \(Q\setminus S\). If \(v\in S\) and \(u\in Q\), then
    \(u\in Q\setminus S\), since \(S\subseteq Y\). Finally, if
    \(v\in S\) and \(u\notin Q\), then \(uv\) is an edge of
    \(H\), so \(u\in N_H(S)\). In every case, \(uv\) has an
    endpoint in \(Q'\). Hence \(Q'\) is a vertex cover smaller
    than \(Q\), as desired.
    \end{proof}


\end{document}